\uuid{CzOQ}
\titre{Opérations sur des sous-espaces vectoriels}
\auteur{hanani ; mzari}
\organisation{exo7}
\licence{CC BY-NC-SA 4.0}
\datecreate{2020-01-20}
\niveau{L1}
\module{Algèbre}
\chapitre{Espace vectoriel}
\sousChapitre{Définition, sous-espace}
\difficulte{2}
\theme{Espaces vectoriels, Opérations sur des sous-espaces vectoriels}
\exo7qcmid{lille-2.espaces-vectoriels.011}

\contenu{
\begin{question}
\qclassification{106.01}
\qtags{motcle=Opérations sur des sous-espaces vectoriels}
Soit $E$ un $\Rr$-espace vectoriel non nul et $F$ et $G$ deux sous-espaces vectoriels de $E$ tels que $F\nsubseteq G$ et  $G\nsubseteq F$. Quelles sont les assertions vraies ?
\begin{answers}
\good{$F+G =\{x+y\, ; \, x\in F \,\mbox {et}\,  y \in G \}$ est un sous-espace vectoriel de $E$.}
\good{$F\cap G $ est sous-espace vectoriel de $E$.}
\bad{$F\cup G $ est sous-espace vectoriel de $E$.}
\good{$F\times G = \{(x,y) \, ; \, x \in F\,\mbox {et}\, y \in G \}$ est un sous-espace vectoriel de $E\times E$.}
\end{answers}
\begin{explanations} La somme, l'intersection et le produit cartésien de sous-espaces vectoriels est un espace vectoriel. Par contre, la réunion de deux sous-espaces vectoriels n'est un espace vectoriel que si l'un des deux est inclus dans l'autre.
\end{explanations}
\end{question}
}
